\caption{Model-adequacy screening, detection losses, and selected-plane uncertainty in controlled synthetic tests. (a) Fraction of held-out T07 records rejected by the five-check adequacy screen under seven generating mechanisms, with separate matched-normal (open) and planted-positive (filled) conditions (48 records per condition and mechanism). Abbreviations: Gaussian VAR(1), cubic saturation, VAR(2), spatially coloured forcing, common-mode bursts, observation noise, and regime switching. (b) Unscreened $R_A>1$ detections (open grey) versus $R_A>1$ detections that also pass the screen (orange), divided in each case by the \emph{full held-out} T07 cohort ($n=216$ planted-positive and $n=432$ matched-normal records; M2). Vertical lines are 95\% subject-cluster bootstrap intervals, except the zero-detection screened normal point, which uses the descriptive Wilson upper bound. The companion version adds descriptive, lightly outlined $R_\sigma>1$ fractions (diamonds), without separate interval estimates in panel (b). (c) Unconditional coverage of nominal 95\% percentile bootstrap intervals for the corrected $R_A$: an oracle fixed planted plane versus the M2 plane reselected for each bootstrap sample, across three observation regimes. Each symbol summarizes 96 independent subject-recordings drawn from six spectral families; error bars are descriptive Wilson intervals. Inadmissible or nonestimable intervals count as noncoverage. The dotted line is nominal 95\% coverage. The companion version shows corresponding $R_\sigma$ coverage as pale diamonds with Wilson intervals. The two modes use different, explicitly stated population targets: the planted plane for the oracle and the population selected-M2 plane for the selected estimator. (d) Numerical M2 solver sensitivity: each faint circle represents the fraction $p_{\mathrm{hit}}$ of 24 independently initialized single starts within 1\% of a strong multistart objective reference in one of 144 audited records (48 per regime). Black bars indicate the medians and interquartile ranges. This is a numerical reference, not a global-optimum certificate. The screen and uncertainty assessments are specific to the synthetic designs shown; no general empirical detection, causality, or interval-calibration guarantee is implied.}